% GENERATED FILE. Edit canonical lessons, then run npm run generate:books.
% canonical-source-hash: fnv1a64:f6a11c6e9331a351
% canonical-lessons: FR-C412-guichet, FR-C412-tarif, FR-C412-trajet, FR-C412-forfait, FR-C412-valider

\chapter{Au guichet}
\label{ch:fr-au-guichet}
\hlchaptermodality{412}

\begin{chapteropening}
\textbf{By the end of this chapter:} \emph{I can buy a ticket at the counter: the ticket office, the price list, the journey, the flat-rate pass and validating the ticket.}

Complete the last lesson of chapter 412: I can buy a ticket at the counter: the ticket office, the price list, the journey, the flat-rate pass and validating the ticket.
\end{chapteropening}

\section[le guichet]{le guichet — the ticket office, the counter window}

\label{lesson:FR-C412-guichet}



\begin{quote}
Before the new one: say the French for the cyclist, then the French for to note.
\end{quote}

\subsection*{You'll want to know: le guichet}
\textbf{le guichet} — \textquotedblleft{}the ticket office, the counter window\textquotedblright{}.

\begin{quote}
Il y a beaucoup de monde au guichet.
\end{quote}

\begin{quote}
\emph{There are a lot of people at the ticket office.}
\end{quote}

\emph{Le guichet automatique} is a cash machine in some places.

\begin{cousinweb}
\emph{guichet} comes from Old French \emph{guischet} 'small door in a big one'. It is usually traced to a Norse or Germanic word, but the origin is not settled; English \emph{wicket} is a cousin (EtymDB).
\end{cousinweb}

\begin{grammarlens}[title={What you've built}]
One more word for telling, explaining and describing at B1.
\end{grammarlens}

\subsection*{Your turn}
\begin{itemize}
\raggedright
  \item \emph{Say it:} \emph{le guichet}
  \item \emph{Say it:} \emph{le guichet}, once more
  \item \emph{Say it:} a sentence of your own with \emph{le guichet}
  \item \emph{Recall:} say the French for the cyclist, then the French for to note, then say \emph{le guichet} again
\end{itemize}

\begin{tcolorbox}[breakable,colback=teal!4,colframe=teal!35!black,title={Before you move on}]
What does le guichet mean? (\textquotedblleft{}The ticket office, the counter window\textquotedblright{}.) Say it once more.
\end{tcolorbox}
\section[le tarif]{le tarif — the fare, the price (list)}

\label{lesson:FR-C412-tarif}



\begin{quote}
Before the new one: say the French for to note, then the French for the ticket office.
\end{quote}

\subsection*{You'll want to know: le tarif}
\textbf{le tarif} — \textquotedblleft{}the fare, the price (list)\textquotedblright{}.

\begin{quote}
Les étudiants paient un tarif réduit.
\end{quote}

\begin{quote}
\emph{Students pay a reduced fare.}
\end{quote}

\emph{Plein tarif} — full price.

\begin{cousinweb}
\emph{tarif} was borrowed from Italian \emph{tariffa}, which came from Arabic \emph{taʿrīfa} 'notice, list of rates' (EtymDB).
\end{cousinweb}

\begin{grammarlens}[title={What you've built}]
One more word for telling, explaining and describing at B1.
\end{grammarlens}

\subsection*{Your turn}
\begin{itemize}
\raggedright
  \item \emph{Say it:} \emph{le tarif}
  \item \emph{Say it:} \emph{le tarif}, once more
  \item \emph{Say it:} a sentence of your own with \emph{le tarif}
  \item \emph{Recall:} say the French for to note, then the French for the ticket office, then say \emph{le tarif} again
\end{itemize}

\begin{tcolorbox}[breakable,colback=teal!4,colframe=teal!35!black,title={Before you move on}]
What does le tarif mean? (\textquotedblleft{}The fare, the price (list)\textquotedblright{}.) Say it once more.
\end{tcolorbox}
\section[le trajet]{le trajet — the ride, the journey (from A to B)}

\label{lesson:FR-C412-trajet}



\begin{quote}
Before the new one: say the French for the ticket office, then the French for the fare.
\end{quote}

\subsection*{You'll want to know: le trajet}
\textbf{le trajet} — \textquotedblleft{}the ride, the journey (from A to B)\textquotedblright{}.

\begin{quote}
Le trajet en train dure deux heures.
\end{quote}

\begin{quote}
\emph{The train journey takes two hours.}
\end{quote}

\emph{Le trajet} is the trip itself; \emph{le voyage} is the whole travel.

\begin{cousinweb}
\emph{trajet} was borrowed from Italian \emph{tragetto} 'crossing', which goes back to Latin \emph{trāiectus} 'a crossing over', from \emph{trāicere} 'to throw across', built on \emph{iacere} 'to throw' (EtymDB).
\end{cousinweb}

\begin{grammarlens}[title={What you've built}]
One more word for telling, explaining and describing at B1.
\end{grammarlens}

\subsection*{Your turn}
\begin{itemize}
\raggedright
  \item \emph{Say it:} \emph{le trajet}
  \item \emph{Say it:} \emph{le trajet}, once more
  \item \emph{Say it:} a sentence of your own with \emph{le trajet}
  \item \emph{Recall:} say the French for the ticket office, then the French for the fare, then say \emph{le trajet} again
\end{itemize}

\begin{tcolorbox}[breakable,colback=teal!4,colframe=teal!35!black,title={Before you move on}]
What does le trajet mean? (\textquotedblleft{}The ride, the journey (from A to B)\textquotedblright{}.) Say it once more.
\end{tcolorbox}
\section[le forfait]{le forfait — the flat-rate pass, the package}

\label{lesson:FR-C412-forfait}



\begin{quote}
Before the new one: say the French for the fare, then the French for the ride.
\end{quote}

\subsection*{You'll want to know: le forfait}
\textbf{le forfait} — \textquotedblleft{}the flat-rate pass, the package\textquotedblright{}.

\begin{quote}
Elle a acheté un forfait pour la semaine.
\end{quote}

\begin{quote}
\emph{She bought a pass for the week.}
\end{quote}

\emph{Un forfait de ski} — a ski pass; \emph{un forfait mobile} — a phone plan.

\begin{cousinweb}
The origin of \emph{forfait} in this sense is not settled. It is usually explained as an old \emph{fur fait} 'done at the agreed rate', reshaped after the older \emph{forfait} 'misdeed'.
\end{cousinweb}

\begin{grammarlens}[title={What you've built}]
One more word for telling, explaining and describing at B1.
\end{grammarlens}

\subsection*{Your turn}
\begin{itemize}
\raggedright
  \item \emph{Say it:} \emph{le forfait}
  \item \emph{Say it:} \emph{le forfait}, once more
  \item \emph{Say it:} a sentence of your own with \emph{le forfait}
  \item \emph{Recall:} say the French for the fare, then the French for the ride, then say \emph{le forfait} again
\end{itemize}

\begin{tcolorbox}[breakable,colback=teal!4,colframe=teal!35!black,title={Before you move on}]
What does le forfait mean? (\textquotedblleft{}The flat-rate pass, the package\textquotedblright{}.) Say it once more.
\end{tcolorbox}
\section[valider]{valider — to validate (a ticket); to approve}

\label{lesson:FR-C412-valider}



\begin{quote}
Before the new one: say the French for the ride, then the French for the flat-rate pass.
\end{quote}

\subsection*{You'll want to know: valider}
\textbf{valider} — \textquotedblleft{}to validate (a ticket); to approve\textquotedblright{}.

\begin{quote}
Dans le bus, les passagers valident leur billet.
\end{quote}

\begin{quote}
\emph{On the bus the passengers validate their tickets.}
\end{quote}

A ticket that is not validated can mean a fine.

\begin{cousinweb}
\emph{valider} is a learned borrowing from Medieval Latin \emph{validāre} 'to make valid', from \emph{validus} 'strong', built on \emph{valēre} 'to be strong' (Lewis \& Short).
\end{cousinweb}

\begin{grammarlens}[title={What you've built}]
One more word for telling, explaining and describing at B1.
\end{grammarlens}

\subsection*{Your turn}
\begin{itemize}
\raggedright
  \item \emph{Say it:} \emph{valider}
  \item \emph{Say it:} \emph{valider}, once more
  \item \emph{Say it:} a sentence of your own with \emph{valider}
  \item \emph{Recall:} say the French for the ride, then the French for the flat-rate pass, then say \emph{valider} again
\end{itemize}

\begin{tcolorbox}[breakable,colback=teal!4,colframe=teal!35!black,title={Before you move on}]
What does valider mean? (\textquotedblleft{}To validate (a ticket); to approve\textquotedblright{}.) Say it once more.
\end{tcolorbox}
